\documentclass[11pt]{article}
\usepackage[a4paper,margin=1in]{geometry}
\usepackage{amsmath,amssymb}
\usepackage{hyperref}
\title{PyTex Validation Program}
\author{PyTex Contributors}
\date{\today}

\begin{document}
\maketitle

\section{Validation Philosophy}

Scientific software should not rely on informal trust. PyTex validation is structured around layered evidence:

\begin{itemize}
  \item unit tests for invariants,
  \item numerical regression tests,
  \item MTEX parity tracking,
  \item interoperability checks,
  \item documentation and figure integrity checks.
\end{itemize}

\section{MTEX Baseline}

Relevant public MTEX tests and examples form the minimum baseline for overlapping functionality. PyTex records this traceability in the parity matrix maintained in the repository documentation.

\section{Beyond MTEX}

PyTex must also validate areas that are central to its own architecture:

\begin{itemize}
  \item provenance retention,
  \item explicit frame normalization,
  \item adapter boundary correctness,
  \item documentation and asset completeness.
\end{itemize}

\end{document}
